\documentclass{standalone}

\usepackage{tikz}
\usepackage{tikz-3dplot}

\usepackage{amsmath, amsfonts, amssymb}

% Tikz
% CUSTOM LENGTHS


\usetikzlibrary{backgrounds,external,shapes.geometric, arrows,calc,positioning,decorations.pathmorphing, scopes}

% PGFPLOTS




\begin{document}
\tdplotsetmaincoords{90}{140}
\begin{tikzpicture}[tdplot_main_coords,scale=2.5]
    \tikzstyle{line} = [thick]
    \tikzstyle{arrow} = [->,>=stealth]
    \tikzstyle{vec} = [thick,->,>={latex[scale=2]}]
    \tikzstyle{lrarrow} = [<->,>=to]
    
    % Sommets des triangles et origine
    \coordinate (O) at (0,0,0);
    \coordinate (a1) at (0,-.5,1.2);
    \coordinate (a2) at (0,.5,1.2);
    \coordinate (a3p) at (1,0,0.2);
    \coordinate (a3m) at (-1,0,0.2);
    
    % Coordonnées pour placer les légendes
    \coordinate (rhonp) at (1.5,0,.7);
    \coordinate (rhonm) at (-1.5,0,.7);
    \coordinate (Tnp) at (1.5,0,1.1);
    \coordinate (Tnm) at (-1.5,0,1.1);
    
    % Labels de l'origine et du vecteur position
    \node[anchor=north] at (O) {$O$};
    % \node[above=0.5 of O] {$x$};
    % \node[anchor=south] at (O) {$x$};
    \node (prout) [above=0.1cm of O] {$x$};

    % Placement des extremites de la longueur ln et du label
    \node (ln1) [above=0.1 of a1] {};
    \node (ln2) [above=0.1 of a2] {};
    \node[anchor=south] (ln) at ($(ln1)!0.5!(ln2)$) {$l_n$};
    
    % coordonnes des barycentres des deux triangles
    \coordinate (bp) at (barycentric cs:a1=1,a2=1,a3p=1);
    \coordinate (bm) at (barycentric cs:a1=1,a2=1,a3m=1);

    % Coordonnees pour les fleches des legendes de Tn
    % Un peu au dessus du  barycentre parce que c'est plus joli
    \coordinate (btp) at ($(bp)!0.1!(ln)$);
    \coordinate (btm) at ($(bm)!0.1!(ln)$);

    % Coordonnees pour l'extremites des vecteurs rn
    % Un peu en dessous du barycentre parce que c'est feng shui
    \coordinate (brp) at ($(bp)!0.2!(a3p)$);
    \coordinate (brm) at ($(bm)!0.2!(a3m)$);

    % Coordonnees pour les fleches des legendes de rho_n
    % Au milieu des rn
    \coordinate (rnp) at ($(brp)!0.5!(a3p)$);
    \coordinate (rnm) at ($(brm)!0.5!(a3m)$);

    % ecriture des labels
    \node[anchor=south] at (rhonp) {$\boldsymbol{\rho}_n^+ = x_n^+-x$};
    \node[anchor=south] at (rhonm) {$\boldsymbol{\rho}_n^- = x-x_n^-$};

    \node[anchor=east] at (Tnp) {$T_n^+$};
    \node[anchor=west] at (Tnm) {$T_n^-$};
    
    \node[anchor=north] at (a3p) {$x_{n^+}$};
    \node[anchor=north] at (a3m) {$x_{n^-}$};
    
    % Dessing des triangles
    % ATTENTION le triangle de derriere doit etre dessine en premier,
    % sinon il apparaît devant
    \draw[fill=white] (a1) -- (a2) -- (a3m) -- cycle;
    \draw[fill=white] (a1) -- (a2) -- (a3p) -- cycle;
    
    % Double fleche ln
    \draw[lrarrow] (ln1) -- (ln2);

    % Vecteurs positions
    \draw[dashed,vec] (O) -- (brp);
    \draw[dashed,vec] (O) -- (brm);

    % Vecteurs rn
    \draw[vec] (brp) -- (a3p);
    \draw[vec] (a3m) -- (brm);

    % Fleches des legendes
    \draw[arrow] (rhonp) to[bend right] (rnp);
    \draw[arrow] (rhonm) to[bend left] (rnm);

    \draw[arrow] (Tnp) to[bend left] (btp);
    \draw[arrow] (Tnm) to[bend right] (btm);
\end{tikzpicture}
\end{document}